\begin{textsample}{NEG Dim 6 – human – Score -3.00 – mg/mg\_ef\_matematica\_3\_p7.txt}  \label{ex:f6_neg_006}
Um dos principais objetivos do ensino de Matemática, em qualquer nível, é o de desenvolver habilidades para a solução de problemas.
Esses problemas podem advir de diferentes situações que \textbf{exijam} o domínio da linguagem matemática e da construção de argumentos que permitam ao estudante elaborar propostas concretas a partir dos conhecimentos adquiridos ao longo do ensino fundamental.
No \textbf{primeiro} caso, é necessária uma boa competência de uso da linguagem matemática para interpretar questões formuladas verbalmente.
No segundo caso, quer dizer que, problemas interessantes que despertam a curiosidade dos estudantes, podem surgir dentro do próprio contexto matemático quando novas situações podem ser exploradas e o conhecimento aprofundado, num exercício contínuo de imaginação e de investigação.

% matched lemmas: exigir, primeiro
\end{textsample}
